\subsection{CLOSED SUBGROUPS}

THEOREM 2. Let G be a finite-dimensional Lie group over R or \(\mathbf{Q}_{p}\) . Every closed subgroup of G is a Lie subgroup of G. More generally, let U be a symmetric open neighbourhood of \(e\) in G and H a non-empty closed subspace of U such that the conditions \(x \in \mathrm{H}\) , \(y \in \mathrm{H}\) and \(xy^{-1} \in \mathrm{U}\) imply \(xy^{-1} \in \mathrm{H}\) . Then H is a Lie subgroup germ of G.

Let \(\mathfrak{h}\) be the Lie subalgebra tangent to H at \(e\) (§ 4, no. 5, Definition 2). There exists a Lie subgroup germ \(\mathrm{H}_0\) of G with Lie algebra \(\mathfrak{h}\) and contained in H. We prove that \(\mathrm{H}_0\) is open in H with the topology induced by that on G. This will prove that H is an analytic submanifold of G and will establish the theorem.

There exist a vector subspace \(\mathfrak{l}\) supplementary to \(\mathfrak{h}\) in \(\mathrm{L}(\mathrm{G})\) , symmetric open neighbourhoods \(\mathrm{V}_1, \mathrm{V}_2\) of zero in \(\mathfrak{h}\) and \(\mathfrak{l}\) respectively and an exponential mapping \(\phi\) of \(\mathrm{G}\) defined on \(\mathrm{V}_1 + \mathrm{V}_2\) and possessing the following properties:

\begin{enumerate}
\def\labelenumi{(\alph{enumi})}
\item
  the mapping \((a_{1}, a_{2}) \mapsto \phi(a_{1})\phi(a_{2})\) is an analytic isomorphism of \(\mathbf{V}_{1} \times \mathbf{V}_{2}\) onto an open subset \(\mathbf{V}\) of \(\mathbf{G}\) ;
\item
  \(\phi (\mathrm{V}_1)\subset \mathrm{H}_0;\)
\item
  \(V^{2} \subset U.\)
\end{enumerate}

We shall prove (and this will achieve the proof) that there exists an open neighbourhood \(\mathbf{V}_2'\) of 0 in \(\mathbf{V}_2\) such that \(\mathbf{H} \cap (\phi(\mathbf{V}_1)\phi(\mathbf{V}_2')) = \phi(\mathbf{V}_1)\) .

Suppose that the assertion is false. Then we can find a sequence \((x_{n})\) in \(\mathbf{V}_1\) and a sequence \((y_{n})\) in \(\mathbf{V}_2 - \{0\}\) tending to 0, such that \(\phi (x_n)\phi (y_n)\in \mathbf{H}\) for all \(n\) . Then \(\phi (y_n)\in \mathbf{H}\) by (c).

If \(\mathrm{K} = \mathbf{Q}_p\) , it can be assumed that \(\mathrm{V}_2\) is an additive subgroup of \(\mathfrak{l}\) and that \(\phi (pa) = \phi (a)^p\) for all \(a\in \mathrm{V}_2\) and all \(p\in \mathbf{Z}\) . Then \(\phi (\lambda y_1)\in \mathbf{H}\) for all \(\lambda \in \mathbf{Z}\) and hence by continuity for all \(\lambda \in \mathbf{Z}_p\) . The mapping \(f\colon \lambda \mapsto \phi (\lambda y_1)\) of \(\mathbf{Z}_p\) into \(G\) is analytic and takes its values in H and \((T_0f)(1) = y_1\) . Hence \(y_{1}\in \mathfrak{h}\) , which is absurd. The theorem is therefore established in the case of \(\mathbf{Q}_p\) .

If \(\mathbf{K} = \mathbf{R}\) , it can be assumed that \(\mathrm{V}_2\) is connected and that \(y_{n}\) belongs to \(\frac{1}{4}\mathrm{V}_2 - \{0\}\) . By taking a subsequence of \((y_n)\) if necessary, we can find a sequence \((\lambda_n)\) of non-zero scalars such that \(\lambda_n^{-1}y_n\) tends to an element \(y\) of \(\mathrm{V}_2 - \{0\}\) . The sequence \((\lambda_n)\) tends to 0. Let \(\lambda \in \mathbf{R}\) be such that \(\lambda y \in \frac{1}{4}\mathrm{V}_2\) ; we prove that \(\exp (\lambda y) \in H\) . It can be assumed that \(\lambda\lambda_n^{-1}y_n \in \frac{1}{4}\mathrm{V}_2\) for all \(n\) . Let \(k_n \in \mathbf{Z}\) be such that \(|\lambda - k_n\lambda_n|\) tends to 0. For \(n\) sufficiently large, \((\lambda - k_n\lambda_n)\lambda_n^{-1}y_n \in \frac{1}{4}\mathrm{V}_2\) and hence \(k_n y_n \in \frac{1}{4}\mathrm{V}_2\) . Hence \(\exp (hy_n) \in H\) for \(h\) an integer and \(0 \leqslant |h| \leqslant |k_n|\) (as is seen by induction on \(|h|\) ). Then

\[
\begin{array}{r l} \exp (\lambda y) & = \lim _ {n \to \infty} \exp (\lambda \lambda_ {n} ^ {- 1} y _ {n}) = \lim _ {n \to \infty} (\exp ((\lambda - k _ {n} \lambda_ {n}) \lambda_ {n} ^ {- 1} y _ {n}) \exp (k _ {n} y _ {n})) \\ & = \lim _ {n \to \infty} \exp k _ {n} y _ {n} \in H. \end{array}
\]

Hence the mapping \(f \colon \lambda \mapsto \exp \lambda y\) , where \(\lambda y \in \frac{1}{4} \mathrm{V}_2\) , takes its values in \(\mathbf{H}\) and \((\mathrm{T}_0f)(1) = y\) . Hence \(y \in \mathfrak{h}\) , which is absurd. The theorem is thus established in the case of \(\mathbf{R}\) .

Theorem 2 is no longer true if \(G\) is not assumed to be finite-dimensional (Exercise 12).

COROLLARY 1. Let G' be a locally compact group, G a finite-dimensional Lie group over R (resp. \(\mathbf{Q}_{p}\) ) and f a continuous morphism of G' into G. If the kernel of f is discrete, G' is a finite-dimensional real (resp. p-adic) Lie group.

There exists a compact neighbourhood V of e in G' such that f\textbar V is a homeomorphism of V onto a compact subspace of G. If U is a sufficiently small open neighbourhood of e in G, the hypotheses of Theorem 2 are satisfied with H = f(V) ∩ U. Hence H is a Lie subgroup germ of G. Let W be the inverse image of H under f \textbar{} V. Then W is a neighbourhood of e in G'. Let W be given the analytic manifold structure transported from that on H by (f \textbar{} W) \(^{-1}\) . For all \(z \in G'\) , the mapping \(x \mapsto f(z)xf(z)^{-1}\) of G into G is analytic; hence there exists an open neighbourhood W' of e in W such that the mapping \(x' \mapsto zx'z^{-1}\) of W' into W is analytic. By Proposition 18 of § 1, no. 9, there exists on G' a Lie group structure which induces on a sufficiently small open neighbourhood of e the same analytic structure as W and hence the same topology as the initial topology on G'.

COROLLARY 2. Let G be a finite-dimensional Lie group over K, H a subgroup of G, V an open neighbourhood of e in G and \((\mathbf{M}_{i})_{i\in I}\) a family of analytic manifolds over K; for all \(i\in I\) , let \(f_{i}\) be a K-analytic mapping of V into \(M_{i}\) such that

\[
\mathrm{H} \cap \mathrm{V} = \{x \in \mathrm{V} | f _ {i} (x) = f _ {i} (e) \text {   for   all   } i \in \mathrm{I} \}.
\]

\begin{enumerate}
\def\labelenumi{(\roman{enumi})}
\item
  If \(\mathbf{K} = \mathbf{C}\) , \(\mathbf{H}\) is a Lie subgroup of \(\mathbf{G}\) .
\item
  If \(\mathbf{K}\) is a finite extension of \(\mathbf{Q}_p\) and \(\mathbf{I}\) is finite, \(\mathbf{H}\) is a Lie subgroup of \(\mathbf{G}\) .
\item
  Suppose that K = C. We consider G as a real Lie group. Then H is a real Lie subgroup of G (Theorem 2). Let \(a \in \mathbf{L}(\mathbf{H})\) . There exists a connected open neighbourhood W of 0 in C such that \(\exp \lambda a \in V\) for all \(\lambda \in W\) . Let \(i \in I\) . Then \(f_i(\exp \lambda a) = f_i(e)\) if \(\lambda \in R \cap W\) . Hence \(f_i(\exp \lambda a) = f_i(e)\) if \(\lambda \in W\) by analytic continuation. Thus, \(\exp \lambda a \in H\) for \(\lambda \in W\) and therefore \(\mu a \in L(H)\) for all \(\mu \in C\) . Therefore H is a Lie subgroup of the complex Lie group G (§ 4, no. 2, Proposition 2).
\item
  Suppose that K is a finite extension of \(Q_{p}\) . We consider G as a Lie group over \(Q_{p}\) . It is finite-dimensional and Theorem 2 implies that H is a p-adic Lie subgroup of G. Since I is finite, \(\prod_{i\in I} M_{i}\) is a manifold and it can be assumed that the family \((f_{i})\) reduces to a single mapping f.~Let \(a \in L(G)\) . Let \(\phi\) be an exponential mapping of G. Then \(f(\phi(\lambda a)) = f(e)\) for \(\lambda \in Q_{p}\) and \(|\lambda|\) sufficiently small. Since \(f\) is K-analytic, it follows that \(f(\phi (\lambda a)) = f(e)\) for \(\lambda \in \mathbf{K}\) and \(|\lambda|\) sufficiently small. Hence \(\phi (\lambda a) \in \mathrm{H}\) for \(\lambda \in \mathrm{K}\) and \(|\lambda|\) sufficiently small and therefore \(\mu a \in \mathrm{L(H)}\) for all \(\mu \in \mathrm{K}\) . The proof is completed as in (i).
\end{enumerate}

Corollary 2 (ii) is no longer true if we omit the hypothesis that I is finite.

